%************************************************
\chapter{Disentanglement}\label{ch:disentanglement}
%************************************************

\textbf{Disentanglement Representation Learning} (DRL) is a learning paradigm in which models are designed to extract and disentangle the 
hidden factors of variation in the data, $\mathbf{z}\in\mathbb{R}^c$, given an input observation $\mathbf{x}\in\mathbb{R}^d$. 
These generative factors are often called latent variables. The formal definition of DRL proposed in \cite{RepresentationLearning} 
is that "Disentangled representation should separate the distinct, independent and informative generative factors of variation in 
the data. Single latent variables are sensitive to changes in single underlying generative factors, while being relatively invariant 
to changes in other factors". This organisation of the latent space into explainable semantic representations improves the following 
properties of a generative model: 
\begin{itemize}
    \item \textbf{Explanability}: the meaningful and independent representations are aligned semantically with latent generative factors. 
    \item \textbf{Generalizability}: the representations are well separated from the original entangled input, which improves the generalization ability for the generation of new samples. 
    \item \textbf{Controllability}: DRL allows for controllable generation by manipulating the learned disentangled representations of the latent space.  
\end{itemize}
There are several disentanglement methods in the current literature, which can be grouped based on the generative model they use.\cite{wang2024disentangledrepresentationlearning}    


\section{VAE-Based Models}  
In Section -- we defined the Variational Autoencoder (VAE) model, where the loss function is defined in Equation --. This type of model has the potential to learn disentangled representations in simple datasets, adding some extra regularisations to the original loss function. \\
One of the first methods was $\beta$\textbf{-VAE} \cite{higgins2017betavae}, which introduces a $\beta$ penalty hyperparameter before the KL term in the ELBO loss. The new training objective is defined as 
\begin{equation}
    \mathcal{L}_{ELBO}(\theta, \phi; \mathbf{x}) = \mathbb{E}_{\mathbf{z}\sim q_{\phi}(\mathbf{z}|\mathbf{x})}\left[\log p_{\theta}(\mathbf{x}|\mathbf{z}) - \beta D_{KL}(q_{\phi}(\mathbf{z}|\mathbf{x})||p(\mathbf{z}))\right],
\end{equation}
where a value of $\beta=1$ corresponds to the original VAE implementation and larger values of $\beta>1$ encourage learning more disentangled representations, adding more pressure to the latent bottleneck to be more factorised, while reducing the performance of reconstruction. This shows an important trade-off between reconstruction accuracy and the disentanglement quality of the latent representation. \\
We can gain deep insight into this trade-off by breaking down the KL term of the loss function as in --
\begin{equation}
    \mathbb{E}_{p_{data}(\mathbf{x})}\left[ KL(q_{\phi}(\mathbf{z}|\mathbf{x})||p(\mathbf{z})) \right] = I(\mathbf{x};\mathbf{z}) + KL(q_{\phi}(\mathbf{z})||p(\mathbf{z})),
\end{equation}
where $I(\mathbf{x};\mathbf{z})$ is the mutual information between $\mathbf{x}$ and $\mathbf{z}$ under the joint distribution $p_{data}(\mathbf{x})q_{\phi}(\mathbf{z}|\mathbf{x})$. Here we can clearly see that making $\beta$ larger pushes $q(\mathbf{z})$ towards the factorial prior $p(\mathbf{x})$, improving the disentanglement, but reduces the amount of information about $\mathbf{x}$ stored in $\mathbf{z}$, producing a poor reconstruction. \\ 
\\
This trade-off between reconstruction and disentanglement is the main motivation for the development of new methods based on VAEs and other generative models. \\
Based on the previous decomposition of the KL term, \textbf{FactorVAE} \cite{disen} proposed a method that only penalises the components that encourage independence in the latent space. The loss function in this method is defined as:
\begin{align}
    \mathcal{L}_{ELBO}(\theta,\phi;\mathbf{x}) &= \mathbb{E}_{\mathbf{z}\sim q_{\phi}(\mathbf{z}|\mathbf{x})}\left[\log p_{\theta}(\mathbf{x}|\mathbf{z}) - D_{KL}(q_{\phi}(\mathbf{z}|\mathbf{x})||p(\mathbf{z}))\right] \\
    &- \gamma D_{KL}(q_{\phi}(\mathbf{z}||\bar q_{\phi}(\mathbf{z}))),
\end{align}
where the last component is the Total Correlation (TC) and $\bar q_{\phi}(\mathbf{z}) = \prod_j q_{\phi}(z_j)$. TC measures the dependencies between the variables in the latent representation. Penalising this specific term encourages the latent dimensions to be statistically independent, which is the core mathematical interpretation of disentanglement. \\
Because the Total Correlation is computationally intractable to evaluate directly over the entire dataset, FactorVAE employs an adversarial training approach. A discriminator network is trained to distinguish between samples drawn from the aggregated posterior $q_{\phi}(\mathbf{z})$ and samples drawn from the product of its marginals $\prod_j q_{\phi}(z_j)$, effectively estimating and minimizing the density ratio. This approach achieves better disentanglement than $\beta$-VAE while preserving a higher reconstruction quality.


\section{Diffusion/Flow-Based Models}  
Diffusion and Flow models are a more recent and advanced class of generative models that have drawn a lot of attention in recent years because of their remarkable performance. Disentangled representation learning can also be applied to these types of models. \\
\\
One of the first approaches was DisDiff \cite{yang2023disdiffunsuperviseddisentanglementdiffusion}, which learns a disentangled representation and a disentangled gradient field for each generative factor. Assume the data samples are generated by $N$ underlying ground truth factors $f^c$, $c\in\{1, \dots, N\}$. Each factor $c$ follows the distribution $p(f^c)$. We can modify the original score function (Equation --) with the conditioning of the factors using the Bayes' Rule as follows: $$\nabla_{\mathbf{x}_t}\log p(\mathbf{x}_t|f^c) = \nabla_{\mathbf{x}_t}\log p(f^c|\mathbf{x}_t) + \nabla_{\mathbf{x}_t}\log p(\mathbf{x}_t),$$ where we already have the score $\nabla_{\mathbf{x}_t}\log p(\mathbf{x}_t)$ from a pretrained unconditional model and we only have to learn $\nabla_{\mathbf{x}_t}\log p(f^c|\mathbf{x}_t)$. Since we are in an unsupervised setting, we do not know $f^c$. We have to learn a set of representations $\{z^c|c=1,\dots,N\}$ that must encompass all information of the original samples $\mathbf{x}_0$ and represent the true factors of variations $f^c$. To obtain these representations we use a set of encoders with learnable parameters $\phi$, defined as $E_{\phi}(\mathbf{x}_0)=\{E_{\phi}^1(\mathbf{x}_0), \dots E_{\phi}^N(\mathbf{x}_0)\}=\{z^1,\dots,z^N\}$. Then, assume we have a factor subset $\mathcal{S}\subseteq\mathcal{C}$ with $z^{\mathcal{S}}=\{z^c|c\in\mathcal{S}\}$, the gradient field can be defined as $$\nabla_{\mathbf{x}_t}\log p(z^{\mathcal{S}}|\mathbf{x}_t) = \sum_{c\in\mathcal{S}}\nabla_{\mathbf{x}_t}\log p(z^c|\mathbf{x}_t).$$ Finally, to estimate the gradient fields, they used a network $G_{\psi}(\mathbf{x}_t, z^c, t)$ and the reconstruction loss is defined as $$\mathcal{L}_r = \mathbb{E}_{\mathbf{x}_0,t,\epsilon} ||\epsilon - \epsilon_{\theta}(\mathbf{x}_t, t) + \frac{\sqrt{\alpha_t}\sqrt{1-\bar\alpha_t}}{\beta_t}\sigma_t\sum_{c\in\mathcal{C}}G_{\psi}(\mathbf{x}_t,z^c,t)||.$$ This loss formulation alone does not guarantee disentanglement among the representations $z^c$. The authors proposed a disentanglement loss that minimizes the upper bound for the mutual information between $z^c$ and $z^j$, where $c\neq j$, defined as $$\mathcal{L}_{dis}=\min_{i,j,\mathbf{x}_0,\mathbf{\hat x}_0^j} ||\hat z^{c|j}-\hat z^c|| - ||\hat z^{c|j} - z^c|| + ||\hat z^{j|j} - z^j||,$$ where $\mathbf{\hat x}_0^j$ is conditioned on $z^j$, $\hat z^c = E_{\phi}^c(\mathbf{\hat x}_0)$ and $\hat z^{c|j} = E_{\phi}^c(\mathbf{\hat x^j_0})$ where $\mathbf{\hat x}_0$ is a sampled data from the pretrained model. \\
\\
EncDiff \cite{yang2024diffusionmodelcrossattention} is an improvement of the previous method where an image encoder $\tau_{\phi}$ aims to provide a set of concept tokens $\mathcal{S}=\{s_1,\dots,s_N\}$. Each dimension of the encoded feature vectors is treated as a disentangled factor and maps each to a vector using a non-shared MLP layer. Then, based on Latent Diffusion Models \cite{rombach2022highresolutionimagesynthesislatent}, the idea is to use the cross-attention mechanism to map the tokens to the intermediate representations of the U-Net model. The general framework of this method is illustrated in Figure \ref{fig:encdiff}. \\
\\
DisenBooth \cite{chen2024disenboothidentitypreservingdisentangledtuning} \\
\\
Flow Matching-Based disentanglement method \cite{chi2026disentangledrepresentationlearningflow}


\section{Metrics} 


%*****************************************
%*****************************************
%*****************************************
%*****************************************
%*****************************************